\section{Introduction}
\label{sec:tpc48-introduction}

The large-prime descent in the TPC sequence produces the determinant
equation
\begin{equation}
 pr-mj=h_0,
 \label{eq:tpc48-intro-determinant}
\end{equation}
with \(m\asymp M\), an orbit length \(J\), and a literal residual
coordinate \(s=d(m)r\).  TPC--46 completed the consecutive-shift
family and found a mesoscopic shift-spectral gap when the output vector
is independent of the orbit coordinate \(j\) \citep{WangTPC46}.
TPC--47 then projected the actual \(j\)-dependent residual coefficient
onto one complete residue-comb band.  Its canonical low core vanishes
there on a nonempty shift interval, while the full frozen family has
arbitrary-power leakage \citep{WangTPC47}.  That theorem is genuinely
sectorial: it neither controls the orthogonal orbit-frequency
complement nor gives a lower bound for the mass retained by the chosen
band.

The natural continuation is to tile the entire orbit-frequency circle
and retain orthogonality through determinant dilation.  Two different
questions are hidden in that sentence.
\begin{enumerate}[label=\textup{(\roman*)},leftmargin=2.4em]
 \item Can the full orbit field be decomposed into exact input-output
       resonance channels without paying the number of bands?
 \item Does that decomposition by itself reduce coherent summation over
       the source slopes \(m\)?
\end{enumerate}
The answer to the first question is yes.  The answer to the second is
no for arbitrary orbit inputs, but becomes quantitatively positive for
the common smooth profiles supplied by the actual projective
reassembly.  This positive/negative distinction is the main result of
the paper.

\subsection{Critical orbit tiling}

Let \(G=H_0^{-1}\Z/\Z\subset\T\), and partition the quotient
\(\T/G\) into \(L\) half-open cells.  Their inverse images
\(\cC_k\), \(k\in\Z/L\Z\), form translated full residue combs of
measure \(1/L\).  If \(\Pi_k\) is the corresponding orbit Fourier
projector, then
\begin{equation}
 \Pi_k\Pi_\ell=0\quad(k\ne\ell),
 \qquad
 \sum_{k\bmod L}\Pi_k=I
 \label{eq:tpc48-intro-resolution}
\end{equation}
in the strong operator topology.  On a scalar orbit window of \(N\)
integers,
\begin{equation}
 \operatorname{tr}(E_{I,v}^*\Pi_kE_{I,v})=\frac NL.
 \label{eq:tpc48-intro-band-trace}
\end{equation}
Here \(E_{I,v}\) denotes zero extension on one fixed unit
non-orbit fiber; it is defined precisely in
\cref{sec:tpc48-tiling}.
The critical choice \(L\asymp J\) therefore has bounded trace per
band.  It sacrifices the growing single-band trace of TPC--47, but no
mass hypothesis is then needed because the orthogonal sum of all bands
is the entire orbit space.

For the modulation
\begin{equation}
 (U_{m,\alpha}y)_j=\e(mj\alpha)y_j,
 \label{eq:tpc48-intro-modulation}
\end{equation}
the block \(\Pi_kU_{m,\alpha}\Pi_\ell\) is supported precisely where
the shifted input cell meets the output cell.  After setting
\(q=k-\ell\pmod L\), the support condition is
\begin{equation}
 d_G(m\alpha,c_q)<\Delta,
 \qquad
 c_q=\frac{q}{LH_0},
 \qquad
 \Delta=\frac1{LH_0}.
 \label{eq:tpc48-intro-tube}
\end{equation}
Here \(d_G(x,y):=\operatorname{dist}(x-y,G)\).
Thus every channel is a union of exact determinant resonance tubes.
At \(m\asymp M\), their half-width in \(\alpha\) is
\begin{equation}
 \frac{\Delta}{m}\asymp\frac1{LM}.
 \label{eq:tpc48-intro-tube-width}
\end{equation}
The critical choice \(L\asymp J\) restores the fine scale
\((JM)^{-1}=X^{-1+o(1)}\).  This is a square-function resolution, not a
spectral gap.

\subsection{The arbitrary-input barrier}

Write \(K=|\cM|\) for the number of supported slopes and define the
unlabelled synthesis
\begin{equation}
 T_\alpha((y_m)_{m\in\cM})
 :=\sum_{m\in\cM}U_{m,\alpha}y_m.
 \label{eq:tpc48-intro-unrestricted-synthesis}
\end{equation}
As an operator from \(\bigoplus_m\widetilde\Hh\) to
\(\widetilde\Hh\), it satisfies the exact identity
\begin{equation}
 \boxed{\|T_\alpha\|^2=K.}
 \label{eq:tpc48-intro-unrestricted-norm}
\end{equation}
Indeed, the upper bound is Cauchy--Schwarz, while
\(y_m=U_{m,\alpha}^{-1}g\) makes every output equal to the same vector
\(g\).  Decomposing each \(y_m\) into all input bands does not change
this example: the input band can compensate the determinant
translation.  Consequently all-band Parseval is not, by itself, a
large-sieve gain.

We make the obstruction more precise by retaining the difference
channel \(q\).  The squared norm of the resulting channel-lift operator
is exactly the essential maximum \(R_L(\alpha)\) of the pointwise
resonance-tube multiplicity.  The physical channel collapse has norm
squared \(L\), and
\begin{equation}
 K\le L R_L(\alpha).
 \label{eq:tpc48-intro-factor-barrier}
\end{equation}
Both constants are optimal in their declared arbitrary-input classes.
No statement here says that the actual TPC coefficient saturates these
classes.

\subsection{The common-profile large sieve}

The inherited projective representation has additional layerwise
structure.  Every projective layer has one common sampled profile
\(\phi_\omega(j/J)\) for all source slopes.  After grouping the prime
and cofactor variables at fixed \(m\), its full shift transform is
\begin{equation}
 \widehat{\cF}^{\omega}(\alpha)
 =\sum_{\substack{j\in\Z\\j\equiv a_\omega\ (H_0)}}
 \phi_\omega(j/J)
 \left(\sum_m B_{\omega,m}(\alpha)\e(mj\alpha)\right)
 \otimes e_j.
 \label{eq:tpc48-intro-layer-normal-form}
\end{equation}
Let \(t=\|\alpha\|_\T\), with \(D_J(0;K):=K\).  Under the no-wrap
condition
\begin{equation}
 H_0(M_+-M_-)t\le\frac14,
 \label{eq:tpc48-intro-no-wrap}
\end{equation}
we prove
\begin{equation}
 \boxed{
 \|\widehat{\cF}^{\omega}(\alpha)\|^2
 \ll_{\phi_\omega,H_0}
 \frac{J}{H_0}D_J(t;K)
 \sum_m\|B_{\omega,m}(\alpha)\|^2,}
 \qquad
 D_J(t;K):=\min\left\{K,1+\frac1{Jt}\right\}\quad(t>0).
 \label{eq:tpc48-intro-large-sieve}
\end{equation}
The proof is a progression Poisson estimate followed by the Schur test
for the vector-valued Gram matrix, in the classical large-sieve
tradition \citep{IwaniecKowalski2004}.  A cluster of consecutive slopes
shows that \(D_J(t;K)\) is sharp up to constants for common-profile
slope boxes containing a sufficiently long consecutive cluster.

The full signed actual coefficient is not one projective layer.
Minkowski reassembles all layers and leaves a declared projective
\(m\)-square envelope.  This is an exact reduction, not an estimate of
that envelope by the physical atomic diagonal.  At fixed \(m\),
different cofactors \(r\) occupy the distinct residual coordinates
\(d(m)r\), so the remaining inner object is an exact \(r\)-square sum
of prime exponential sums.

\subsection{Endpoint ledger and claim boundary}

At the inherited endpoint,
\begin{equation}
 M=X^{267/400+o(1)},
 \qquad
 J=X^{133/400+o(1)}.
 \label{eq:tpc48-intro-endpoint}
\end{equation}
The critical band count is \(L=X^{133/400+o(1)}\).  On a coarse-edge
shift scale \(t\asymp M^{-1}\), a dense slope box has normalized
collision factor
\begin{equation}
 D_J(t;K)=X^{67/200+o(1)}.
 \label{eq:tpc48-intro-critical-load}
\end{equation}
The TPC allowance is only \(X^{1/400+o(1)}\), leaving the exact exponent
gap
\begin{equation}
 \frac{67}{200}-\frac1{400}=\frac{133}{400}.
 \label{eq:tpc48-intro-allowance-gap}
\end{equation}
Thus, in the dense sharpness class, the large sieve gains the full
orbit length relative to arbitrary slopewise Cauchy, but it does not
close atomic normalization.  This ledger is not a lower bound for the
literal slope support, whose cardinality may be smaller.

The analysis concerns the frozen shift family built from the literal
physical \(h_0\)-coefficient field.  Only \(h=h_0\) is the inherited
physical residual synthesis, and that object belongs to the
low-sign-averaged residual Gram route rather than the native affine
physical face.  Nothing below proves an atomic Bessel bound, controls
the ultra-low or alias complements, reconstructs physical packets at
new shifts, breaks the parity barrier, or establishes a prime-pair
asymptotic.
